\subsubsection{Selección de tecnologías del sistema}

Se aplica el método de dos filtros de la metodología de desarrollo. Los criterios de eliminación derivan de los requerimientos del sistema (Tabla~\ref{tab:requerimientos}), y los pesos de los criterios de ponderación, de las restricciones del proyecto. Se evalúan por separado el framework del servidor, la aplicación del navegador y el motor de base de datos; las tecnologías del agente conversacional se seleccionan en el apartado siguiente. Cada criterio indica entre paréntesis su origen, y en todas las matrices 3 es la puntuación más favorable.

\paragraph{Framework del servidor}

Se comparan Spring Boot, Django, Laravel y NestJS (Tabla~\ref{tab:ob2_matriz_backend}). La experiencia del desarrollador recibe el mayor peso, porque un único desarrollador debe cumplir el plazo; la agenda, el núcleo compartido por la recepción y el agente, recibe el segundo.

\begingroup
\fontsize{9}{11}\selectfont
\renewcommand{\arraystretch}{1.2}
\setlength{\LTpre}{0pt}\setlength{\LTpost}{6pt}
\setlength{\tabcolsep}{4pt}
\begin{xltabular}{\linewidth}{|>{\raggedright\arraybackslash}X|>{\centering\arraybackslash}p{1.2cm}|>{\centering\arraybackslash}p{1.8cm}|>{\centering\arraybackslash}p{1.8cm}|>{\centering\arraybackslash}p{1.8cm}|>{\centering\arraybackslash}p{1.8cm}|}
\caption{Evaluación de frameworks de servidor mediante dos filtros}
\label{tab:ob2_matriz_backend} \\
\hline
\textbf{Criterio} & \textbf{Peso} & \textbf{Spring Boot} & \textbf{Django} & \textbf{Laravel} & \textbf{NestJS} \\ \hline
\endfirsthead
\hline
\textbf{Criterio} & \textbf{Peso} & \textbf{Spring Boot} & \textbf{Django} & \textbf{Laravel} & \textbf{NestJS} \\ \hline
\endhead
\multicolumn{6}{|l|}{\cellcolor{gray!10}\textbf{Criterios de eliminación (primer filtro)}} \\ \hline
Control de acceso por roles y autenticación por token mediante componentes maduros (RNF01) & --- & Cumple & Cumple & Cumple & Cumple \\ \hline
Código abierto sin licencias de pago (restricción de costo) & --- & Cumple & Cumple & Cumple & Cumple \\ \hline
\multicolumn{6}{|l|}{\cellcolor{gray!10}\textbf{Criterios de ponderación (segundo filtro)}} \\ \hline
Transacciones y validaciones de dominio robustas para la agenda (RF07, RF11) & 20\,\% & 3 & 3 & 2 & 2 \\ \hline
Auditoría de cambios integrable (RNF02) & 15\,\% & 3 & 3 & 3 & 2 \\ \hline
Experiencia previa del desarrollador (restricción de un único desarrollador) & 25\,\% & 3 & 2 & 1 & 2 \\ \hline
Consumo de memoria en el servidor (restricción de costo) & 15\,\% & 1 & 2 & 3 & 2 \\ \hline
Pruebas de integración con base de datos real (RF07: restricción anticruce) & 10\,\% & 3 & 3 & 2 & 2 \\ \hline
Lenguaje compartido con el frontend en TypeScript (restricción de un único desarrollador) & 15\,\% & 1 & 1 & 1 & 3 \\ \hline
\textbf{Total ponderado (1--3)} & \textbf{100\,\%} & \textbf{2,40} & \textbf{2,30} & \textbf{1,90} & \textbf{2,15} \\ \hline
\end{xltabular}
\endgroup

Spring Boot obtiene el mayor total, con el consumo de memoria más alto del grupo. La diferencia con Django es estrecha y la decide la experiencia del desarrollador: sin ese criterio, Django obtendría el mayor total. Ese criterio es, por tanto, el decisivo, y se justifica porque un único desarrollador debe cumplir el plazo. Laravel gana en consumo de memoria y NestJS en lenguaje compartido con el frontend.

\paragraph{Aplicación del navegador}

Se comparan Angular, React y Vue (Tabla~\ref{tab:ob2_matriz_frontend}). Los formularios del Formulario 033 y la experiencia del desarrollador reciben los mayores pesos.

\begingroup
\fontsize{9}{11}\selectfont
\renewcommand{\arraystretch}{1.2}
\setlength{\LTpre}{0pt}\setlength{\LTpost}{6pt}
\setlength{\tabcolsep}{4pt}
\begin{xltabular}{\linewidth}{|>{\raggedright\arraybackslash}X|>{\centering\arraybackslash}p{1.2cm}|>{\centering\arraybackslash}p{1.8cm}|>{\centering\arraybackslash}p{1.8cm}|>{\centering\arraybackslash}p{1.8cm}|}
\caption{Evaluación de frameworks de frontend mediante dos filtros}
\label{tab:ob2_matriz_frontend} \\
\hline
\textbf{Criterio} & \textbf{Peso} & \textbf{Angular} & \textbf{React} & \textbf{Vue} \\ \hline
\endfirsthead
\hline
\textbf{Criterio} & \textbf{Peso} & \textbf{Angular} & \textbf{React} & \textbf{Vue} \\ \hline
\endhead
\multicolumn{5}{|l|}{\cellcolor{gray!10}\textbf{Criterios de eliminación (primer filtro)}} \\ \hline
Aplicación de una sola página en el navegador, sin instalación en los equipos de la clínica & --- & Cumple & Cumple & Cumple \\ \hline
Código abierto sin licencias de pago (restricción de costo) & --- & Cumple & Cumple & Cumple \\ \hline
\multicolumn{5}{|l|}{\cellcolor{gray!10}\textbf{Criterios de ponderación (segundo filtro)}} \\ \hline
Formularios extensos con validación integrada, como el Formulario 033 (RF01, RNF06) & 25\,\% & 3 & 1 & 2 \\ \hline
TypeScript nativo para tipar los datos de la API (restricción de un único desarrollador) & 15\,\% & 3 & 2 & 2 \\ \hline
Experiencia previa del desarrollador (restricción de un único desarrollador) & 25\,\% & 3 & 2 & 2 \\ \hline
Curva de aprendizaje (restricción de plazo) & 15\,\% & 1 & 2 & 3 \\ \hline
Componentes disponibles para calendario y odontograma (RF02, RF07) & 20\,\% & 2 & 3 & 2 \\ \hline
\textbf{Total ponderado (1--3)} & \textbf{100\,\%} & \textbf{2,50} & \textbf{1,95} & \textbf{2,15} \\ \hline
\end{xltabular}
\endgroup

Angular obtiene el mayor total por sus formularios integrados y su soporte nativo de TypeScript, aunque su curva de aprendizaje es la más exigente. React gana en componentes disponibles, y Vue, en la curva de aprendizaje.

\paragraph{Motor de base de datos}

Se comparan PostgreSQL, MySQL y MongoDB (Tabla~\ref{tab:ob2_matriz_bd}). La garantía de no cruce recibe el mayor peso, porque la recepción y el agente reservan sobre la misma agenda.

\begingroup
\fontsize{9}{11}\selectfont
\renewcommand{\arraystretch}{1.2}
\setlength{\LTpre}{0pt}\setlength{\LTpost}{6pt}
\setlength{\tabcolsep}{4pt}
\begin{xltabular}{\linewidth}{|>{\raggedright\arraybackslash}X|>{\centering\arraybackslash}p{1.2cm}|>{\centering\arraybackslash}p{2.0cm}|>{\centering\arraybackslash}p{2.0cm}|>{\centering\arraybackslash}p{2.0cm}|}
\caption{Evaluación de motores de base de datos mediante dos filtros}
\label{tab:ob2_matriz_bd} \\
\hline
\textbf{Criterio} & \textbf{Peso} & \textbf{PostgreSQL} & \textbf{MySQL} & \textbf{MongoDB} \\ \hline
\endfirsthead
\hline
\textbf{Criterio} & \textbf{Peso} & \textbf{PostgreSQL} & \textbf{MySQL} & \textbf{MongoDB} \\ \hline
\endhead
\multicolumn{5}{|l|}{\cellcolor{gray!10}\textbf{Criterios de eliminación (primer filtro)}} \\ \hline
Integridad referencial declarativa (claves foráneas) para datos clínicos estructurados (RF01, BN1) & --- & Cumple & Cumple & No cumple \\ \hline
Código abierto sin licencias de pago (restricción de costo) & --- & Cumple & Cumple & Cumple \\ \hline
\multicolumn{5}{|l|}{\cellcolor{gray!10}\textbf{Criterios de ponderación (segundo filtro)}} \\ \hline
Garantía de no cruce ante reservas concurrentes, en la propia base de datos (RF07) & 35\,\% & 3 & 1 & --- \\ \hline
Búsqueda por similitud de texto para el catálogo CIE-10 (RF03) & 25\,\% & 3 & 2 & --- \\ \hline
Consumo de recursos en un servidor compartido (restricción de costo) & 20\,\% & 2 & 3 & --- \\ \hline
Herramientas de respaldo y exportación en formato abierto (RNF04) & 20\,\% & 3 & 3 & --- \\ \hline
\textbf{Total ponderado (1--3)} & \textbf{100\,\%} & \textbf{2,80} & \textbf{2,05} & \textbf{---} \\ \hline
\end{xltabular}
\endgroup

PostgreSQL obtiene el mayor total. Es la única alternativa con una restricción de exclusión nativa, necesaria para impedir cruces de citas en la propia base de datos, y su extensión de similitud de texto facilita la búsqueda en el catálogo CIE-10. MySQL gana en consumo de recursos. MongoDB queda descartado en el primer filtro, porque no ofrece integridad referencial declarativa.

La Tabla~\ref{tab:ob2_stack} resume la pila tecnológica seleccionada.

\begingroup
\fontsize{9}{11}\selectfont
\renewcommand{\arraystretch}{1.2}
\setlength{\LTpre}{0pt}\setlength{\LTpost}{6pt}
\setlength{\tabcolsep}{4pt}
\begin{xltabular}{\linewidth}{|p{2.3cm}|p{3.4cm}|X|p{2.6cm}|}
\caption{Pila tecnológica seleccionada para el sistema}
\label{tab:ob2_stack} \\
\hline
\textbf{Capa} & \textbf{Tecnología seleccionada} & \textbf{Componentes clave} & \textbf{Requerimientos atendidos} \\ \hline
\endfirsthead
\hline
\textbf{Capa} & \textbf{Tecnología seleccionada} & \textbf{Componentes clave} & \textbf{Requerimientos atendidos} \\ \hline
\endhead
Backend & Spring Boot 3.5 con Java 21 & Spring Security con JWT, Spring Data JPA, validación de datos, migraciones con Flyway e historial de cambios con Hibernate Envers (usuario y fecha de cada revisión) & RF07, RF11, RNF01, RNF02 \\ \hline
Frontend & Angular 22.1 & Enrutamiento, formularios y servicios con inyección de dependencias; aplicación de una sola página servida por nginx & RF02, RF07, RNF06 \\ \hline
Base de datos & PostgreSQL 16 & Restricción de exclusión anticruce (extensión \texttt{btree\_gist}); búsqueda por similitud con \texttt{pg\_trgm} e índices GIN; migraciones con Flyway & RF01, RF03, RF07, RNF04 \\ \hline
Contenedores & Docker y Docker Compose & Todos los servicios en un único archivo de composición, portable a otro proveedor & RNF04, RNF05 \\ \hline
Pruebas & JUnit 5 y Testcontainers & Pruebas de integración contra PostgreSQL real, incluida la restricción anticruce de la agenda & RF07, RF11 \\ \hline
\end{xltabular}
\endgroup
